% =============================================================
%  Chapter 8 — Conclusion
% =============================================================

\chapter{Conclusion}
\label{chap:conclusion}

% -----------------------------------------------------------------------
\section{Revisiting the Thesis Problem}
\label{sec:ch8:problem}
% -----------------------------------------------------------------------

Autonomous agents operating in unknown environments cannot rely on
fixed policies, fixed rewards, or fixed planners.
They must act before the full structure of the world is known, update
their behavior as geometry and risk are revealed, and do so without
destabilizing the decisions already learned.
This requires not merely a learned mapping from observations to actions,
but a \emph{continually correctable decision dynamics}: a vector field
over the agent's state space whose topology can be reshaped as new
information arrives.

Existing paradigms provide important pieces of this capability.
Imitation learning offers sample efficiency; reinforcement learning
offers long-horizon optimization; classical planning offers explicit
structure; safety-critical control offers formal invariance guarantees;
Hamiltonian and port-Hamiltonian learning offers structured,
passivity-regularized dynamics.
But no single paradigm provides all four simultaneously: structured
dynamics, continual adaptation, passivity-monitored updates, and
risk-sensitive objectives.

We studied that gap through the lens of
\textbf{continual Hamiltonian learning}: decision-making modeled as a
structured phase-space dynamics whose energy landscape is progressively
reshaped by geometry, material properties, and risk.
The policy is the flow map of a reduced Hamiltonian; adapting the
policy means reshaping the energy landscape, not relearning a flat
action mapping.

% -----------------------------------------------------------------------
\section{Summary of Contributions}
\label{sec:ch8:contributions}
% -----------------------------------------------------------------------

We developed this viewpoint through two stages, connected by a
formal enrichment principle, and validated empirically.

\paragraph{Stage 1 — Geometry-only structured navigation (Chapter~\ref{chap:grlsnam}).}
GRL-SNAM established that navigation can be cast as learning a reduced
Hamiltonian $H = H_{\mathrm{kin}} + H_{\mathrm{geom}}$ under local
sensing alone.
The policy is the symplectic flow map; a three-module architecture
coordinates sensing, frame motion, and shape control at different
timescales; online geometric alignment updates barrier weights from
rollout feedback without catastrophic relearning.
Empirically, this geometry-only learner outperforms reactive planners
and deep RL baselines by a large margin under matched sensing
constraints, and does so at fraction of the workspace coverage
required by global planners.

\paragraph{Stage 2 — Material-aware risk-sensitive learning (Chapter~\ref{chap:material}).}
The richer learner enlarged the stored energy with soft-risk and
hard-hazard barrier terms, the objective with a CVaR tail-risk
criterion, the coefficient space with material force heads, and the
adaptation hierarchy with a three-loop curriculum.
Every addition followed the same structural chain: a new energy term
inducing a new force channel, a new coefficient, a new objective term,
a new gradient pathway, and a new update law.
Empirically, this reduces cumulative risk exposure by 56.4\% over the
geometry-only baseline and $7\times$ the oscillation count, while
path length increases by only 1.1\%.

\paragraph{The enrichment principle (Chapter~\ref{chap:hierarchy}).}
Theorem~\ref{thm:ch5:enrichment} formalized the progression:
any new energy term forces five downstream consequences, making each
transition between stages a structural necessity rather than an
implementation choice.
Backward-compatible initialization guarantees that passivity is
preserved at the start of each enrichment, and the passivity-guided
curriculum monitors stability throughout.

\paragraph{Experimental validation (Chapter~\ref{chap:experiments}).}
The ablation study confirmed that every added structural component
resolves a specific failure mode soft-risk energy removes hazard
drift, CVaR removes catastrophic tail encounters, the passivity
criterion prevents unstable phase advancement, and the explicit
dissipative port eliminates oscillation.
The one-to-one correspondence between structure and failure mode is
the strongest empirical evidence that the hierarchy is not superficial.

\paragraph{Research agenda (Chapter~\ref{chap:limitations}).}
Five theoretical open problems passivity-preserving update guarantees,
multi-timescale convergence rates, basin-of-competence volume bounds,
force-decomposition identifiability, and Sobolev field regularity 
frame the next stages of the research program.

% -----------------------------------------------------------------------
\section{Central Lessons and Conceptual Takeaways}
\label{sec:ch8:lessons}
% -----------------------------------------------------------------------

\paragraph{Lesson 1: The right object to learn is the field, not the policy.}
Our primary reframing is that the policy should be identified
with the flow map of a structured energy landscape, not with an
unstructured action distribution.
This identification makes the force-field topology explicit, exposes
stability through passivity analysis, and provides a principled handle
for continual adaptation: reshaping the energy landscape rather than
retraining a black box.

\paragraph{Lesson 2: Added complexity should be added structurally.}
We demonstrate that when the environment becomes harder when
geometry-only sensing is insufficient for safe navigation the correct
response is not to add a new module on top of an existing system, but
to \emph{enlarge the learner} along four coupled axes simultaneously:
stored energy, coefficient space, objective, and adaptation timescale.
Enlarging one axis without the others produces the failure modes
documented in the ablation study: hazard drift without the material
energy term, catastrophic tail encounters without CVaR, instability
without the passivity criterion.

\paragraph{Lesson 3: The progression itself is the contribution.}
The individual navigation numbers—98.7\% success on geometry, 56.4\%
risk reduction on material—are evidence, not the claim.
The claim is the enrichment principle: that a structured phase-space
learner can be progressively enlarged in a way that is principled,
backward-compatible, and passivity-regularized (empirically stability-monitored).
That principle is what makes the framework a \emph{research program}
rather than a pair of navigation systems.

\paragraph{Lesson 4: Passivity is a practical stability language.}
The port-Hamiltonian separation of stored energy and dissipative port,
combined with the passivity surrogate $\mathcal{L}_{\mathrm{pass}}$,
provided an actionable stability monitor throughout training.
It is not a full Lyapunov certificate that remains an open problem 
but it is precise enough to guide curriculum advancement, detect
instability early, and prevent the phase transitions that most
destabilize continual learning.

% -----------------------------------------------------------------------
\section{Final Outlook}
\label{sec:ch8:outlook}
% -----------------------------------------------------------------------

We have shown that structured Hamiltonian learning provides a
principled route into the unknown-environment navigation problem.
It does not claim to solve that problem fully.
Oracle material access, simple observer dynamics, isotropic geometry,
and empirical coverage surrogates are honest limitations.
But the framework's organizing principle that each new source of
environmental complexity should be matched by a structural enlargement
of the energy landscape, the coefficient space, the objective, and the
update hierarchy is general enough to extend beyond these limitations.

Stage~3 of the ladder (perception-led material inference), Stage~4
(belief-state observer and anisotropic metric), and the theoretical
program of convergence guarantees and identifiability analysis are
not distant aspirations.
They are the next applications of the same enrichment chain that
connected Stage~1 to Stage~2.
The structure of what needs to be done is visible because the
enrichment principle makes it visible.

\begin{quote}\itshape
We studied navigation in unknown environments through the
lens of continual Hamiltonian learning.
Rather than treating decision-making as the one-shot learning of a
fixed policy or reward, we developed a structured phase-space
viewpoint in which geometry, safety, and semantic risk act by
reshaping an energy-driven vector field.
Beginning from a geometry-only foundation and extending to a richer
material-aware risk-sensitive learner, we showed how added
complexity in the environment should be matched by added structure
in the stored energy, the objective, and the update laws.
The broader contribution is not only a stronger navigation method,
but a principled framework for progressively enriching decision
dynamics while preserving interpretability, stability, and
continual adaptability and a formal account of why that progression
is structurally necessary rather than incidentally convenient.
\end{quote}